\section{اشتقاق مصفوفات النموذج الديناميكي }
\label{sec:lagrange_derivation}

بناءً على صياغة لاغرانج-أويلر التي تم تقديمها في القسم الأول، يتطلب حساب القوى والعزوم المحركة تفكيك تابع لاغرانج وحساب مشتقاته الجزئية بالنسبة لمتغيرات المفاصل وسرعاتها.  يمكن كتابة العزم الكلي الفعال المطبق على المفصل $(i)$ قبل دمج المحركات على النحو التالي:

\begin{equation}
\tau_i = \sum_{j=1}^{n} D_{ij}\ddot{q}_j + \sum_{j=1}^{n} \sum_{k=1}^{n} D_{ijk}\dot{q}_j\dot{q}_k + D_i
\label{eq:full_analytical_lagrange}
\end{equation}


\subsection{عناصر مصفوفة العطالة $(D_{ij})$}
تمثل المعاملات $(D_{ij})$ عناصر مصفوفة القصور الذاتي الكلية للروبوت $M(q)$ بأبعاد $3 \times 3$. تُعبر هذه العناصر عن عطالة المحاور والتداخل القصور الذاتي المتبادل بين المفاصل أثناء التسارع. تُحسب هذه المعاملات تحليلياً بالاعتماد على أثر مصفوفة التحويل المتكاملة (\LR{Trace}) ومصفوفات الاشتقاق الحركي $(U_{pi})$ وفق القانون التالي:

\begin{equation}
D_{ij} = \sum_{p=\max(i,j)}^{n} \text{Trace} \left( U_{pj} J_p U_{pi}^T \right)
\label{eq:inertia_element_law}
\end{equation}

حيث يمثل $J_p$ مصفوفة القصور الذاتي المتكاملة للوصلة $(p)$ والتي تم تحديدها في القسم السابق، بينما تمثل المصفوفة $U_{pi}$ المشتق الجزئي لمصفوفة التحويل الكلية بالنسبة للمتغير المفصلي $q_i$، والتي تعكس مدى تأثر تموضع الوصلة بتغير زاوية أو إزاحة المفصل.

\subsection{معاملات قوى كوريوليس والطرد المركزي $(D_{ijk})$}
تُعبر المعاملات ثلاثية الدليل $(D_{ijk})$ عن تأثيرات قوى كوريوليس (عند $j \neq k$) وقوى الطرد المركزي (عند $j = k$) الناتجة عن التداخل الحركي للوصلات الدورانية أثناء الدوران السريع. تُحسب هذه المعاملات قياسياً عبر العلاقة:

\begin{equation}
D_{ijk} = \sum_{p=\max(i,j,k)}^{n} \text{Trace} \left( U_{pjk} J_p U_{pi}^T \right)
\label{eq:coriolis_element_law}
\end{equation}

وكما تم إقراره في الفرضيات الأساسية لهذا الفصل، ونظراً لتشغيل الروبوت ضمن مجالات سرعة منخفضة إلى متوسطة، فقد تم إهمال هذه الحدود بالكامل في هذا المشروع، مما يعني هندسياً أن:
\begin{equation}
\sum_{j=1}^{n} \sum_{k=1}^{n} D_{ijk}\dot{q}_j\dot{q}_k = 0
\end{equation}

\subsection{عناصر شعاع الجاذبية الأرضية $(D_i)$}
تمثل المعاملات $(D_i)$ عناصر شعاع العزوم الناتجة عن الجاذبية الأرضية $G(q)$ المؤثرة على مفاصل الذراع. بالاعتماد على متجهات مراكز الثقل المحددة سابقاً $r_p$ لكل وصلة، واعتبار الأرض هي المرجعية الصفرية للارتفاع، تم استنتاج عناصر شعاع الجاذبية تحليلياً وفق القانون التالي:

\begin{equation}
D_i = \sum_{p=i}^{n} -m_p g^T U_{pi} r_p
\label{eq:gravity_element_law}
\end{equation}

حيث أن $g = [0, 0, -g]^T$ هو شعاع تسارع الجاذبية الأرضية، و $m_p$ هي كتلة الوصلة (مع أخذ الكتلة المركبة المتغيرة $m_3^*$ للوصلة الثالثة بعين الاعتبار).

\subsection{النمذجة المصفوفية المستنتجة للروبوت}
بتعويض القوانين التحليلية السابقة وتجميع العناصر، يختزل النموذج الديناميكي الصرف لذراع الروبوت (قبل ربطه بالجهود الكهربائية للمحركات) إلى البنية المصفوفية المدمجة التالية:

\begin{equation}
\begin{bmatrix} \tau_1 \\ \tau_2 \\ \tau_3 \end{bmatrix}_{robot} = 
\begin{bmatrix} 
D_{11} & D_{12} & D_{13} \\ 
D_{21} & D_{22} & D_{23} \\ 
D_{31} & D_{32} & D_{33} 
\end{bmatrix}
\begin{bmatrix} \ddot{\theta}_1 \\ \ddot{\theta}_2 \\ \ddot{d}_3 \end{bmatrix} + 
\begin{bmatrix} D_1 \\ D_2 \\ D_3 \end{bmatrix}
\label{eq:robot_matrix_form}
\end{equation}

تُشكل هذه المنظومة الرياضية المقيدة الرابط الفيزيائي الأساسي بين حركة ذراع الروبوت الميكانيكية والعزوم الخارجية الفعالة المطلوبة لتنفيذ المناورة التشغيلية.